% Generated draft: requires mathematical and source review.
\documentclass[11pt]{article}
\usepackage[margin=25mm]{geometry}
\usepackage{fontspec}
\setmainfont{Latin Modern Roman}
\usepackage{amsmath,amssymb,mathtools,unicode-math}
\setmathfont{Latin Modern Math}
\usepackage{enumitem,needspace,xurl,hyperref}
\setlength{\parindent}{0pt}
\setlength{\parskip}{5pt}
\setlength{\emergencystretch}{3em}
\newcommand{\N}{\mathbb{N}}
\newcommand{\Z}{\mathbb{Z}}
\newcommand{\Q}{\mathbb{Q}}
\newcommand{\R}{\mathbb{R}}
\newcommand{\C}{\mathbb{C}}
\newcommand{\F}{\mathbb{F}}
\newcommand{\A}{\mathbb{A}}
\newcommand{\PP}{\mathbb{P}}
\newcommand{\E}{\mathbb{E}}
\newcommand{\eps}{\varepsilon}
\newcommand{\dd}{\,\mathrm{d}}
\begin{document}
\section*{UnsolvedMath ID: 30000451}
\textbf{Source problem code:} catalogue entry
\subsection{Definitions and mathematical statement}
Use the report's algebraic-group highest-weight setting in characteristic $p$, with Coxeter number $h$, dominant weights $X^+$ and $p$-restricted dominant weights
\[
 X_1^+=\{\nu\in X^+:0\le\langle\nu,\alpha^\vee\rangle<p
                              \text{ for every simple coroot }\alpha^\vee\}.
\]
Write $L(\nu)$ for the simple highest-weight module, $\Delta(\nu)$ for the Weyl module and $\nabla(\nu)$ for the dual Weyl module. The superscript $(1)$ denotes Frobenius twist. Regular means $p$-regular in the affine-dot-action convention: the shifted weight avoids the affine root hyperplanes. In standard notation this is $\langle\lambda+\rho,\alpha^\vee\rangle\notin p\Z$ for every root, where $\rho$ is half the sum of positive roots.

Assume $p>h$. For each regular dominant $\lambda=\lambda_0+p\lambda_1$, with $\lambda_0\in X_1^+$ and $\lambda_1\in X^+$, define exactly the prime-marked modules
\[
 \Delta'_{\mathrm{red}}(\lambda)
       =L(\lambda_0)\otimes\Delta(\lambda_1)^{(1)},\qquad
 \nabla'_{\mathrm{red}}(\lambda)
       =L(\lambda_0)\otimes\nabla(\lambda_1)^{(1)}.
\]
Conjecture 3.6 asks whether, for every such weight,
\[
 \Delta'_{\mathrm{red}}(\lambda)[-l(\lambda)]\in E^L,
 \qquad
 \nabla'_{\mathrm{red}}(\lambda)[-l(\lambda)]\in E^R.
\]
Here $[-l(\lambda)]$ is the source's derived shift, $l(\lambda)$ its affine-length convention, and $E^L,E^R$ its specified left/right derived subcategories. Their generators or homological membership criteria are not given in the extract, and cannot be replaced by invented Ext-vanishing definitions. The full group/root-datum hypotheses and block/truncation conventions must likewise be matched to the report.

The original extraction also mentions the distinct Conjecture 3.7 filtration request, separately stated in companion permanent ID $30000452$: $\Delta(\lambda)$ should have a finite filtration with sections $\Delta_{\mathrm{red}}(\nu)$, and $\nabla(\lambda)$ one with sections $\nabla_{\mathrm{red}}(\nu)$, for regular dominant $\nu$. Those are unprimed reduced modules; they are not identified with the tensor-defined prime modules above. Companion ID $30000450$ likewise concerns unprimed-module membership in $E^L,E^R$. The verified primary target of this record is precisely Conjecture 3.6, while those bundled neighboring requests are retained as related context.

The literature assessment checked on 21 August 2026 says the surrounding conditional results do not prove the full regular-weight homological and filtration assertions. No restriction to the Jantzen region or additional character-formula hypothesis is inserted as though it solved the stated all-regular-weight question. The open label is dated triage requiring independent verification.
\subsection{Short English statement}
For $p>h$, do the tensor-defined prime-reduced standard and costandard modules, shifted by minus the prescribed affine length, lie in the report's left/right derived subcategories for every regular dominant weight? Prime and unprimed reduced modules, and the separate filtration conjecture, must remain distinct.
\subsection{Sources}
\begin{itemize}
\item \emph{Representations of Finite Groups} (2007), Oberwolfach Reports, DOI 10.4171/owr/2006/15, pp. 895--978, Conjecture 3.6, \url{https://doi.org/10.4171/owr/2006/15}; proposed year supplied as 2006.
\item Literature assessment checked 21 August 2026: original report, \url{https://doi.org/10.4171/owr/2006/15}; Parshall--Scott, ``Beyond the Jantzen Region,'' \url{https://ems.press/content/serial-article-files/46045}.
\item Companion local records: ID $30000450$, \texttt{OWR-1265-009}, gives the unprimed Conjecture 3.2 membership request; ID $30000452$, \texttt{OWR-1265-012}, gives Conjecture 3.7's regular reduced-module filtrations.
\item Local UnsolvedMath registry, permanent ID $30000451$, exact source code \texttt{catalogue entry}, created 25 August 2026 and updated 16 September 2026. Supplied verification identifies Conjecture 3.6 and confirms the prime marks, weight restrictions, shifts and subcategories.
\end{itemize}
\end{document}
